\documentclass[11pt,letterpaper]{article}
\usepackage[T1]{fontenc}
\usepackage[utf8]{inputenc}
\usepackage{lmodern,microtype}
\usepackage[margin=0.9in]{geometry}
\usepackage{amsmath,amssymb,amsthm,mathtools}
\usepackage{booktabs,longtable,array,enumitem}
\usepackage{xcolor,hyperref,xurl,fancyhdr,listings}
\definecolor{navy}{HTML}{16324F}
\definecolor{teal}{HTML}{087E8B}
\definecolor{light}{HTML}{F3F6F8}
\hypersetup{colorlinks=true,linkcolor=navy,urlcolor=teal,citecolor=navy,pdftitle={Further Lean proof simplification: controller arithmetic and subset coverage},pdfauthor={Research report prepared for Vladimir Reshetnikov}}
\setlength{\parindent}{0pt}
\setlength{\parskip}{6pt}
\setlength{\emergencystretch}{2em}
\setlist{topsep=3pt,itemsep=3pt,leftmargin=*}
\pagestyle{fancy}
\fancyhf{}
\fancyhead[L]{\small\sffamily ProveIt / Further proof simplification}
\fancyhead[R]{\small\sffamily 26 October 2026}
\fancyfoot[C]{\small\thepage}
\setlength{\headheight}{14pt}
\newtheorem{theorem}{Theorem}[section]
\newtheorem{lemma}[theorem]{Lemma}
\newtheorem{corollary}[theorem]{Corollary}
\newcommand{\N}{\mathbb N}
\newcommand{\Z}{\mathbb Z}
\newcommand{\Zm}[1]{\Z/#1\Z}
\newcommand{\card}[1]{\lvert #1\rvert}
\newcommand{\repobase}{https://github.com/VladimirReshetnikov/ProveIt/blob/9f35798db4a63c25e435fca14d75931144657d7e}
\newcommand{\refreshbase}{https://github.com/VladimirReshetnikov/ProveIt/blob/f00ba53b331833c6cb64707bd90e48489e5610f5}
\newcommand{\freshleanfile}[1]{\href{\refreshbase/Computability/DiophantineConsistency/Lean/DiophantineConsistency/#1.lean}{\nolinkurl{#1.lean}}}
\newcommand{\leanfile}[1]{\href{\repobase/Computability/DiophantineConsistency/Lean/DiophantineConsistency/#1.lean}{\nolinkurl{#1.lean}}}
\newcommand{\ident}[1]{\nolinkurl{#1}}
\newcommand{\artifact}[1]{\nolinkurl{#1}}
\lstset{basicstyle=\small\ttfamily,backgroundcolor=\color{light},frame=single,rulecolor=\color{light},breaklines=true,columns=fullflexible,keepspaces=true,showstringspaces=false,xleftmargin=5pt,xrightmargin=5pt,aboveskip=8pt,belowskip=8pt}
\begin{document}
\hypersetup{pageanchor=false}
\begin{titlepage}
\vspace*{0.65in}
{\Large\sffamily\color{teal} PROVEIT RESEARCH NOTE\par}
\vspace{0.25in}
{\Huge\bfseries\color{navy} Further Lean proof\\[6pt] simplification\par}
\vspace{0.2in}
{\Large Controller arithmetic and subset coverage\par}
\vspace{0.35in}
{\large Prepared for Vladimir Reshetnikov\par}
{\large 26 October 2026\par}
\vspace{0.35in}
\begin{minipage}{0.94\textwidth}
\textbf{Main finding.} Allowing one gap in a finite interval gives every
fixed-cardinality subset sum in a cyclic group, for every cardinality
$0<a<T$. Coprimality is unnecessary. Applied to the existing affine power
coordinates, this removes the optional Boolean correction from the arithmetic
coverage theorem, provided its cardinality bound permits $a=C_0$.

\textbf{Other findings.} The fixed CRT components need periods, while the
auxiliary prime-power component needs exact order. A single fixed residue
modulo $94517\cdot120199$ proves the entire two-parameter family of controller
coefficients is coprime to the required moduli. An all-exponent divisibility
law reorganizes the prime-power order proof without a zero/successor split.

\textbf{Review baseline.}
\href{https://github.com/VladimirReshetnikov/ProveIt/commit/9f35798db4a63c25e435fca14d75931144657d7e}{\texttt{9f35798db4a63c25e435fca14d75931144657d7e}}.
The initial audit covered 23 placed packets and five incoming archives.
A final refresh through commit \texttt{f00ba53b331833c6cb64707bd90e48489e5610f5}
reviewed the expanded 36-packet corpus and new controller clients. Earlier
proposals are treated as available even when marked pending or deferred.
\end{minipage}
\vfill
{\small This report proposes proof changes and supplies reproducible Lean
experiments. It does not assert new mathematical priority, a completed
Diophantine consistency equivalence, or a measured whole-project speedup.}
\end{titlepage}
\hypersetup{pageanchor=true}
\pagenumbering{roman}
\begingroup
\small
\setlength{\parskip}{0pt}
\tableofcontents
\endgroup
\newpage
\pagenumbering{arabic}

\section{Result, scope, and priorities}

The most valuable simplification is to strengthen the finite subset theorem
that the controller consumes. The existing proof realizes a sum by translating
a consecutive interval. Translation changes its sum by $a$ times the shift,
so that particular construction needs $a$ invertible modulo $T$. Changing the
shape of the subset by just one gap removes that restriction. Consequently,
the arithmetic coverage proof no longer needs a preliminary Boolean choice
whose purpose is to make the cardinality coprime to an auxiliary prime power.

This is a concrete improvement over the stipulated prior-work baseline.
It does not rename an old scanner, coefficient, witness-transport, or
finite-table abstraction. The new proof applies to cardinalities that the
old construction cannot handle. Sections~\ref{sec:subset} and~\ref{sec:affine}
give the explicit witness and its exact interface consequence.

\begin{longtable}{@{}>{\raggedright\arraybackslash}p{0.14\textwidth}>{\raggedright\arraybackslash}p{0.34\textwidth}>{\raggedright\arraybackslash}p{0.43\textwidth}@{}}
\toprule
Priority & Proposal & What disappears from the relevant proof path \\
\midrule
\endhead
P1 & Arbitrary proper cardinalities in cyclic subset sums & Cardinality coprimality and, with the endpoint change, the optional Boolean correction. \\
P2 & Periods for fixed CRT components & Proofs of minimal multiplicative order where consumers only need an annihilating exponent. \\
P3 & One residue for the complete controller coefficient & Separate dimension, quadratic-factor, coupling-factor, and coprimality arguments at each of three primes. \\
P4 & Uniform annihilator law for the principal-unit tower & The $k=0$ versus successor split in the exact-order proof. \\
Companion & Nonzero first lift as the public seed & Repeated natural-subtraction and divisibility bridges at the current-generator boundary. \\
Optional & Translating a specified seed & Interval-specific triangular arithmetic when a shape constraint makes translation useful. \\
\bottomrule
\end{longtable}

P1--P4 are the primary directions. The companions target overlapping code
and are not counted as additional independent source savings. In particular, the arbitrary-seed translation and unrestricted-cardinality
constructions target the same interval machinery. Their deletion opportunities
are not additive. A prime-support diagnostic suggestion was withdrawn after
the final refresh found it in an earlier submitted packet.

The proof checks are pinned to the original revision; the final novelty and
client review uses the refreshed revision. The Lean project pins
Lean~4.32.0 and mathlib commit
\nolinkurl{81a5d257c8e410db227a6665ed08f64fea08e997} in
\href{\repobase/lean-toolchain}{\nolinkurl{lean-toolchain}} and
\href{\repobase/lake-manifest.json}{\nolinkurl{lake-manifest.json}}.
Validation receipts, source hashes, and the precise replay scope are supplied
with the archive; Section~\ref{sec:validation} explains their interpretation.

\section{Excluding the previous proposals}

\subsection{Final refresh and a withdrawn candidate}

The final freshness check found 25 additional commits, through
\href{https://github.com/VladimirReshetnikov/ProveIt/commit/f00ba53b331833c6cb64707bd90e48489e5610f5}{\texttt{f00ba53b3318}}.
The placed corpus now contains 36 packets. Five of the 13 new placements
are the already-audited incoming archives; the other eight are new reports,
\texttt{zc} through \texttt{zj}. The final exclusion review covers this
delta as well. The exact original inventory below is retained as provenance,
and the additional hashes and comparisons are supplied in the refresh audit.
That delta contains 34 primary TeX components and 219 searched textual
members. The five incoming relocations were checked byte-for-byte for their
retained members; earlier omitted material remains part of the assumed
baseline. See \artifact{evidence/FINAL_REFRESH_NOVELTY_AUDIT.txt} and
\artifact{evidence/refresh_delta_inventory.json} for the final decisions.

The refresh identified a real overlap: \href{\refreshbase/Computability/DiophantineConsistency/research/lean-proof-simplification/2026-10-26zd/report.tex}{packet \texttt{zd}} already proposes
finite prime-support certificates for multiplicative-order diagnostics.
That candidate was therefore removed from this report's recommendations and
from the delivered Lean experiments, despite its successful local check.
All 15 source/configuration inputs used by
the focused replay are unchanged between the two pins. The proof receipts
remain anchored to the original pin; the novelty and client-impact review
uses the refreshed corpus.

The newer \texttt{zh} controller proposal still keeps the Boolean adjustment
and coprime-cardinality path, and its Hs simplification works on the separate
prime residues. These are assumed available here. They do not supply P1's
noncoprime subset theorem or P3's whole-coefficient residue transport.

\subsection{What was audited}

At the initial pin, the supplied exclusion locations contain 23 placed packets, named
\texttt{2026-10-26a} through \texttt{2026-10-26w}, and five additional
proof-simplification ZIPs in \texttt{docs/incoming}. The inventory records
74 primary TeX files or components, their paths, SHA-256 hashes, sizes, and
section headings. Targeted searches cover 1,295 textual members, including
supporting notes and archived source snapshots. That number describes the
search scope; it is not a claim that every word of every duplicated snapshot
was read. Report-section review and mechanism-level comparisons supplement
the searches.

The five incoming files are:
\begin{itemize}
\item \artifact{Lean_Proof_Simplification_Post_Adoption_2026-10-26.zip};
\item \artifact{ProveIt_Additional_Lean_Simplifications_2026-10-26.zip};
\item \artifact{diophantine-lean-further-deltas-2026-10-26.zip};
\item \artifact{diophantine-lean-proof-refinements-2026-10-26.zip};
\item \artifact{proveit-lean-structural-simplifications-2026-10-26.zip}.
\end{itemize}

The machine-readable evidence is in \artifact{evidence/exclusion_inventory.json},
\artifact{evidence/exclusion_taxonomy.json}, and
\artifact{evidence/prior_report_headings.tsv}. The large incoming README also
links unrelated historical research; the audit does not claim to reread every
transitively linked manuscript. The novelty claim is relative to the directly
audited simplification corpus, not to all mathematical literature.

\subsection{How ``assume already applied'' was interpreted}

Every earlier proposed proof interface is treated as available, including
optional, deferred, uncompiled, or unselected ideas. This is a conceptual
union of prior methods. It is not a claim that all competing historical
patches can be combined into one compiling implementation. Accordingly,
this report does not subtract lines from the current tree and label the result
an exact saving against an imaginary fully merged tree.

Excluded families include list and decoder APIs; scanner and history fusion;
primitive-recursive congruence; assignment naturality; common DAG, gate,
residual-square, and degree machinery; coefficient fibers and reflection;
carry and cyclic-word interfaces; Pell sign, minimum, and perturbation
arguments; internal-model arithmetic specifications; finite inverse
invariance; Wang-tile fibers, helical periods, shared edges, pulse rows,
receiver causality, and operating-row recognition. Earlier generic advice to
use finite checks, weaken an unused premise, or introduce an equivalence is
also not presented as a discovery here.

The retained proposals have more specific mechanisms: noncoprime
fixed-cardinality cyclic subset sums; the exact separation between the fixed
and auxiliary order obligations; erasure of both controller exponents in one
residue ring; and a uniform all-exponent divisibility interface for LTE.
The accompanying audit gives the nearest old ideas and the distinction for
each proposal.

\section{P1: all proper cardinalities realize every cyclic sum}
\label{sec:subset}

\subsection{The limitation of translated intervals}

\leanfile{ControllerIntervalSums} constructs a subset by translating
$\{0,1,\ldots,a-1\}$. Its sum is
\[
 a\,s+\sigma_a,\qquad \sigma_a=\sum_{i=0}^{a-1}i=\frac{a(a-1)}2.
\]
Solving for the shift $s\in\Zm T$ requires multiplication by $a$ to be
surjective. The existing theorem therefore uses $\gcd(a,T)=1$ and an
explicit unit inverse. The division by two is correctly performed in the
naturals, not in the residue ring; that existing safeguard is not itself
a new proposal.

The unit condition is a restriction on this family of translated subsets,
not on the existence of a subset with the requested cardinality and sum.

\begin{theorem}[Fixed-cardinality cyclic subset sums]
For integers $T,a$ with $0<a<T$ and any $u\in\Zm T$, there is a subset
$S\subseteq\Zm T$ such that
\[
 \card S=a,\qquad \sum_{x\in S}x=u.
\]
No coprimality, primality, or prime-power hypothesis is required.
\end{theorem}

\subsection{An explicit witness with one gap}

First prove a natural-number statement. For every
$0\le\delta\le a(T-a)$, the integer $\sigma_a+\delta$ is the sum of an
$a$-element subset of $\{0,\ldots,T-1\}$. Write
\[
 \delta=aq+r,\qquad 0\le r<a,
\]
and take
\begin{equation}
 x_i=i+q+\mathbf 1_{\{i\ge a-r\}},\qquad 0\le i<a.
 \label{eq:staircase}
\end{equation}
The sequence is strictly increasing: consecutive entries differ by one,
except that a single jump may have size two. Its sum is
\[
 \sum_{i<a}x_i=\sigma_a+aq+r=\sigma_a+\delta.
\]
The bound on $\delta$ implies $q\le T-a$. If $q=T-a$, then $r=0$;
otherwise $q\le T-a-1$. In the first case the largest entry is $T-1$;
in the second it is at most $(a-1)+(T-a-1)+1=T-1$. Thus all entries are
valid distinct representatives below $T$.

For the modular theorem, take $\delta$ to be the canonical representative of
$u-\sigma_a$ modulo $T$. It is at most $T-1$, and
\begin{equation}
 a(T-a)-(T-1)=(a-1)(T-a-1)\ge0.
 \label{eq:window}
\end{equation}
The natural theorem therefore applies. Casting the chosen representatives
to $\Zm T$ preserves distinctness and gives the desired sum.

This proof is elementary and constructive. It uses a quotient and remainder,
a finite image, an injectivity proof, and a suffix-cardinality sum. It needs
no additive-combinatorics theorem and does not infer a field structure on
\texttt{ZMod T}.

\subsection{A small example that distinguishes the methods}

For $T=4$ and $a=2$, translated consecutive pairs have sums $2s+1$ and
therefore reach only odd residues. Formula~\eqref{eq:staircase} reaches all
four residues:

\begin{center}
\begin{tabular}{@{}ccc@{}}
\toprule
Target $u$ & Selected representatives & Sum modulo $4$ \\
\midrule
$0$ & $\{1,3\}$ & $0$ \\
$1$ & $\{0,1\}$ & $1$ \\
$2$ & $\{0,2\}$ & $2$ \\
$3$ & $\{1,2\}$ & $3$ \\
\bottomrule
\end{tabular}
\end{center}

This is the relevant gain: an even cardinality modulo four becomes usable.
The method does not just shorten an inverse calculation that remains
necessary elsewhere in the same argument.

\subsection{Lean decomposition and exact targets}

\artifact{lean/RestrictedSubsetSums.lean} separates the proof into
\ident{staircase_strictMono}, \ident{staircase_sum},
\ident{staircase_lt}, \ident{exists_nat_subset_sum},
\ident{window_covers_modulus}, and \ident{exists_subset_card_sum}.
The natural theorem permits $a=T$ with offset zero. The modular theorem
correctly requires $a<T$: a subset of size $T$ is the whole group and has
only one possible sum.

The first replacement target is the solver behind
\ident{ControllerIntervalSums.targetInterval}. In
\leanfile{ControllerOrbitTransport}, the affine-sum transport remains useful;
its subset input changes from the inverse-selected interval to the new
staircase subset. The source search finds that the direct non-check client of
\ident{targetInterval} is this orbit-transport module. Historical explicit
interval APIs can be retained as compatibility results if they are still
wanted, but retaining every one reduces the achievable source deletion.

\section{The affine consequence: remove the Boolean arithmetic correction}
\label{sec:affine}

\subsection{The new endpoint}

\leanfile{ControllerResidueLift} currently proves
\ident{exists_subset_boolean_sum}. Its hypotheses include a prime $p$,
$p\mid C_0$, $T=p^n$, an auxiliary weight with nonzero projection modulo
$p$, and an injective coordinate map. Its conclusion uses a Boolean
coefficient $e\le1$ to arrange a cardinality $0<a<C_0$ coprime to $T$.

With the previous section's theorem, the arithmetic endpoint can instead be:

\begin{theorem}[Boolean-free affine coverage]
Suppose $0<C_0<T$, $\iota:\Zm T\to A$ is injective, and weights satisfy
\[
 w(\iota(z))=1+C_0z\quad\text{in }\Zm{(C_0T)}.
\]
Then every $t\in\Zm{(C_0T)}$ is the weight sum of an actual finite subset
$S\subseteq A$ with
\[
 1\le\card S\le C_0.
\]
The expression $C_0z$ denotes the well-defined additive map from $\Zm T$
to $\Zm{(C_0T)}$; it does not require choosing a ring homomorphism between
these residue rings in the wrong direction.
\end{theorem}

\begin{proof}
Let $a$ be the least positive representative of the projection of $t$ to
$\Zm{C_0}$. Thus $1\le a\le C_0$ and $t-a$ has zero projection. There
is a coordinate $u\in\Zm T$ with $t=a+C_0u$. Since $a\le C_0<T$,
the cyclic subset theorem supplies a subset $U\subseteq\Zm T$ of size
$a$ and sum $u$. Set $S=\iota(U)$. Injectivity preserves cardinality, and
additivity gives
\[
 \sum_{x\in S}w(x)=\sum_{z\in U}(1+C_0z)
 =a+C_0\sum_{z\in U}z=a+C_0u=t.
\]
\end{proof}

\subsection{Why the non-strict endpoint is necessary}

The existing output bound $a<C_0$ cannot be preserved while removing the
correction term for all targets. Every weight projects to $1$ modulo $C_0$,
so a subset sum projects to its cardinality. A target with zero projection
requires $a\equiv0\pmod{C_0}$. Its least positive permitted cardinality is
therefore $C_0$ itself.

The checked adapter in \artifact{lean/AffineSubsetCoverage.lean} reuses the
existing \ident{exists_coordinate} and \ident{scaledCoordinate}. The old
coordinate lift uses a representative $a_0<C_0$. When $a_0=0$, the adapter
sets $a=C_0$ and replaces the coordinate $u$ by $u-1$; additivity shows
$C_0+C_0(u-1)=C_0u$. This explicitly handles the endpoint rather than
asking simplification tactics to hide a change in the theorem statement.

If a downstream capacity theorem truly requires $a<C_0$, keep a Boolean
choice in that API. The new subset-sum theorem still removes the need to
make $a$ coprime to $T$, so that weaker migration can eliminate the
prime-power cardinality-selection argument without changing the output
bound. The fully Boolean-free endpoint should be adopted only where the
capacity contract accepts $a\le C_0$.

\subsection{The actual controller parameters and remaining composition}

The existing \leanfile{CurrentPowerCoordinates2011} supplies the relevant
orbit with
\[
 p=120199,\quad d=94517,\quad C_0=dp,\quad T=p^k.
\]
The value of $C_0$ is $dp$, not $d$. Since the repository proves $p>d$,
the inequality $C_0<T$ holds at $k\ge2$. The power-coordinate result also
has smaller-exponent cases, but that does not make this subset theorem
applicable when $T\le C_0$.
The orbit modulus here is $C_0T=dp^{k+1}$. When identifying it with the
coefficient modulus $dp^n$ in P3, set $n=k+1$; the coverage range is then
$n\ge3$. P3's separate unit theorem remains valid for every $n$, including
zero.

For this arithmetic coverage result, the hypotheses $p\mid C_0$, primality,
the prime-power presentation of $T$, and a separately supplied nonzero
Boolean weight disappear. They remain where independently needed to prove
the actual orbit or other controller facts. The new theorem does not show
that every selected coordinate corresponds to a safe physical dummy
position, nor that a specific shape constraint is preserved. It also does
not by itself reduce the variables in the printed Diophantine equation.

The bound $C_0<T$ has content. With $C_0=T=2$, the affine weights are
$\{1,3\}$ modulo four, whose subset sums miss $2$. Thus simply removing
that hypothesis would make the proposed statement false.

\subsection{Compatibility with the newer physical-position proofs}

The refreshed source now contains
\freshleanfile{CurrentControllerSubsets2011}. Its
\ident{exists_normalized_subset} still calls the Boolean-based residue-lift
theorem, so it is a direct new client for the proposed endpoint. The
subsequent normalization by the unit $B^2$ is retained; its separate
$\varepsilon B^3$ summand disappears when $\varepsilon=0$.

There is a useful compatibility fact in
\freshleanfile{CurrentControllerPositions2011}:
\ident{support_safe} and \ident{cyclicSupport_safe} quantify over an
arbitrary finite subset of orbit indices. They require the geometric
condition $r_w<k$, but no interval shape, coprimality, or strict cardinality
cap. Their statements therefore apply to the staircase subset with the
optional parameter set to zero. This is a source-level compatibility
deduction from the maintained signatures, not a claim that the complete
new composition has been rebuilt.

The worst-case physical support size does not increase. The existing
\ident{support_card} theorem gives $a+\varepsilon$. Under the old bounds,
$a<C_0$ and $\varepsilon\le1$ imply $a+\varepsilon\le C_0$; under the
new bounds, $a\le C_0$ and $\varepsilon=0$ give the same cap. Thus allowing
the arithmetic endpoint $a=C_0$ does not by itself demand more physical
updates than the old interface already allowed.

The refresh also includes safe single-dummy increments, nonwrapping
transport, and rank decrements in \freshleanfile{CurrentDummySafety2011}
and \freshleanfile{CurrentDummyMutation2011}. These are existing ingredients
to reuse, not missing local bounds to prove again. The unchecked migration
step for this proposal is the finite-update composition: preserve the
unused-dummy premise at each selected position, aggregate the increments,
and bind that sum to the Boolean-free power-subset endpoint. The late-head
scope addendum records the exact declarations and source identities.

\section{P2: use periods at fixed factors, exact order at the auxiliary factor}

\subsection{The current over-specification}

\leanfile{CurrentControllerData2011} proves the exact orders of $2$ and
$B=2^{94517}$ modulo $47$ and $2011$. The production consumers of
\ident{B_order_47} and \ident{B_order_2011} in the reviewed subtree are
the fixed-prime power reductions in \leanfile{CurrentControllerUnits2011}
and the projection lemmas in \leanfile{CurrentPowerCoordinates2011}.
Those consumers only need
\[
 B^{23}=1\pmod{47},\qquad B^{402}=1\pmod{2011}.
\]
They never use the minimality of $23$ or $402$.

Two short monoid facts supply the appropriate interface:
\[
 x^m=1\land m\mid e\ \Longrightarrow\ x^e=1,
 \qquad
 x^m=1\ \Longrightarrow\ (x^s)^m=1.
\]
The proofs are exponent laws; no group or commutativity assumption is
needed. The file \artifact{lean/FixedPeriods.lean} proves both and the actual
fixed-prime certificates. It also contains an alternative direct bounded
modular certificate for the complete projection modulo $94517\cdot120199$.
That calculation works in the finite residue ring and does not construct
the unreduced natural number $2^{11360754366}$.

\subsection{Where exactness must remain}

\leanfile{ControllerPowerCoordinates} proves injectivity of the exponent
parameter using the exact equation $\operatorname{ord}(g)=T$ and exponent
representatives below $T$. Here an annihilating exponent is insufficient.
For example, $1^3=1$ modulo nine, but the order is one and the three powers
do not define an injection.

The correct division of responsibilities is therefore specific:
\begin{itemize}
\item fixed CRT components: prove the generator projects to $1$;
\item auxiliary prime-power component: prove exact order $p^k$;
\item CRT composition: reuse the existing least-common-multiple order theorem;
\item coordinate enumeration: consume the resulting exact full order.
\end{itemize}

The CRT least-common-multiple theorem is already present as
\ident{CurrentPowerCoordinates2011.order_product}. Recommending it again
would not satisfy the user's request. Likewise, the report does not propose
replacing the auxiliary exact-order premise by a mere period.

There are two distinct implementation goals. Rewriting the consumers removes
their logical dependence on the fixed-prime minimality proofs. To remove
the corresponding build dependency as well, move optional exact-order
diagnostics into a separate module. If they remain in the same imported
source file, Lean still elaborates them when that file builds. No timing
claim is inferred from a shorter theorem dependency chain.

\section{P3: one fixed residue proves the whole controller unit family}

\subsection{Erase both exponents in one residue ring}

Set
\[
 p=120199,\quad d=94517,\quad M=dp=11360848883,\quad
 B=2^{94517},\quad r=2^{45919}.
\]
The last equality matches the maintained factor
$\mathrm{radix}^{977}$ because $\mathrm{radix}=2^{47}$ and
$47\cdot977=45919$. Keep $K$ as the original maintained kernel.

The coefficient in \leanfile{CurrentControllerUnits2011} is
\[
 \Gamma(r_w,n)=r\bigl(q^2-1\bigr)
       \bigl(1+q(K+B^{p^{r_w}})\bigr),\qquad q=B^{p^n}.
\]
The width exponent $r_w$ and the cell-count exponent $n$ are independent;
neither is assumed positive.

In any monoid, a single fixed-point equation $x^p=x$ implies
\[
 x^{p^n}=x\quad\text{for every }n\in\N.
\]
Induct on $n$ and use $(x^{p^n})^p=x^{p^{n+1}}$. No primality, unit,
or finite-field hypothesis is needed for this induction. For the actual
parameters, bounded modular exponentiation verifies
\[
 (B:\Zm M)^p=B.
\]
Both $B^{p^n}$ and $B^{p^{r_w}}$ therefore reduce to $B$ in the same ring,
and the whole expression reduces at once:
\begin{equation}
 \Gamma(r_w,n)\equiv
 \Gamma_0:=r(B^2-1)(1+BK+B^2)\pmod M.
 \label{eq:gamma}
\end{equation}

The natural subtraction is accounted for before casting to the residue
ring: $B\ge1$ implies $B^2\ge1$ and $(B^{p^n})^2\ge1$. Rewriting a
natural subtraction as ring subtraction without those guards would be
invalid. For example, the natural expression $0^2-1$ is zero, whereas the
ring expression is $-1$.

\subsection{One coprimality transport lemma}

\begin{lemma}
If $a\equiv b\pmod{dp}$ and $\gcd(b,dp)=1$, then
\[
 \gcd(a,dp^n)=1\quad\text{for every }n\in\N.
\]
\end{lemma}
\begin{proof}
Congruence preserves the gcd with $dp$, so $a$ is coprime to both $d$ and
$p$. It is then coprime to $p^n$ and to $dp^n$.
\end{proof}

This implication does not require $d$ and $p$ to be coprime or $p$ to be
prime. In the actual application, $M$ is the radical of $dp^n$ for $n>0$.
At $n=0$ the target modulus has smaller prime support, and the same
one-way implication still works. An unconditional converse would be false:
$3$ is coprime to $2\cdot3^0$, but not to $2\cdot3$.

Combining the lemma with~\eqref{eq:gamma} reduces the entire two-parameter
unit theorem to one finite certificate for $\Gamma_0$.

\subsection{The actual finite certificate}

The current kernel residue facts in \leanfile{CurrentKernelResidues2011}
are $K\bmod47=26$, $K\bmod2011=1995$, and
$K\bmod120199=114787$. CRT and finite arithmetic give:

\begin{center}
\begin{tabular}{@{}lr@{}}
\toprule
Quantity modulo $M=11360848883$ & Representative \\
\midrule
$K$ & $1491904576$ \\
$B$ & $10287028452$ \\
$r$ & $9930929229$ \\
$\Gamma_0$ & $706142772$ \\
\bottomrule
\end{tabular}
\end{center}

Finally,
\[
 \gcd(706142772,11360848883)=1.
\]
\artifact{lean/ResidueTower.lean} checks these computations in Lean while
keeping $K$ symbolic and using exactly the three existing residue facts as
premises. It does not replace $K$ by the smaller CRT representative as a
natural number; only its residue is replaced inside modular arithmetic.

The generic theorem \ident{gamma_coprime} and the actual-parameter adapter
\ident{current_gamma_coprime} include both zero-exponent endpoints. The
adapter is a checked standalone result conditional on the existing kernel
facts. A production integration must still apply it to the actual definition
of $K$, identify the definitions of $r$ and $B$, and rebuild the callers.

\subsection{Expected proof simplification}

The replacement targets four families, each with three prime-specific instances:
\ident{power_dimension_*}, \ident{quadratic_mod_*},
\ident{coupling_mod_*}, and \ident{coprime_Gamma_*} lemmas. Their separate
prime-specific derivations become one fixed-point proof, one whole-expression
congruence, and one gcd certificate. Existing diagnostic residue statements
can be recovered as projections if their public names should remain.

This proposal overlaps P2 in the fixed-prime dimension calculations; those
lines must not be counted twice. The gain established here is the simpler
dependency structure and the checked replacement core, not an already
measured global deletion or build-speed improvement. The theorem proves a
coefficient is a unit. It supplies no physical placement, safe increment,
or final arbitrary-target rank correction.

\section{P4: make the all-exponent annihilator law the LTE interface}

\subsection{A uniform theorem}

For an odd prime $p$ and a natural $g$ with
\[
 p\mid g-1,\qquad p^2\nmid g-1,
\]
the maintained \leanfile{PrimePowerOrder} proves exact lifting at exponents
$p^k$, then derives exact order by splitting $k$ into zero and successor.
The proposed interface exposes the whole divisibility predicate first:
\begin{equation}
 p^{k+1}\mid g^m-1\quad\Longleftrightarrow\quad p^k\mid m,
 \qquad k,m\in\N.
 \label{eq:annihilator}
\end{equation}

For positive $m$, LTE says
\[
 v_p(g^m-1)=v_p(g-1)+v_p(m)=1+v_p(m),
\]
and comparing valuations gives~\eqref{eq:annihilator}. The formal proof
uses extended multiplicity so that $m=0$, where both divisibilities hold,
is included uniformly. The natural-subtraction guard still comes from
$g>1$, which follows from $p^2\nmid g-1$.

Let $o=\operatorname{ord}(g\bmod p^{k+1})$. Applying the forward direction
to exponent $o$ gives $p^k\mid o$, and applying the reverse direction to
exponent $p^k$ gives $o\mid p^k$. Antisymmetry proves
$o=p^k$. This avoids the specialized prime-power order criterion and the
$k=0$ case split.

This is not a stronger final mathematical endpoint than the existing
exact-order theorem: the latter already implies~\eqref{eq:annihilator}
through the usual order-divides-exponent characterization. The new
contribution is the organization of the proof from LTE, giving a reusable
all-exponent interface for essentially the same source size. The supplied
refactor preserves the original public statements and names; exact
helper-inclusive size and validation are recorded with the patch.

\subsection{Accept a nonzero first lift directly}

The current auxiliary generator is $g=B^{120199-1}$, and its checked
remainder is
\[
 g\bmod120199^2=7300767062=1+120199\cdot60739.
\]
\artifact{lean/LiftSeed.lean} gives a small generic adapter:
\[
 1<p,\quad c>0,\quad g\bmod p^2=1+pc
 \quad\Longrightarrow\quad
 g>1,\ p\mid g-1,\ p^2\nmid g-1.
\]
The residue lower bound yields $g>1$. Reducing again modulo $p$ gives
the first divisibility. The second would force the remainder modulo $p^2$
to be one, contradicting $pc>0$. A separate $c<p$ premise is unnecessary:
the fact that $1+pc$ is a remainder already supplies that bound.

The target is \leanfile{CurrentAuxiliaryOrder2011}, which currently proves
the positivity and divisibility bridges separately. This adapter aligns
the theorem input with the finite certificate already available. The
generic helper has a real source cost; at its first use alone, a large net
line reduction is not claimed.

Oddness and primality remain necessary for the order corollary. Taking
$p=2,g=3$ gives the same nonzero first-lift pattern, but the order modulo
eight is two instead of four. Taking $p=9,g=28$ likewise shows that an odd
composite base cannot replace the prime assumption.

\section{A companion interface when a specified subset shape matters}

P1 permits the subset shape to change. If an application instead needs a
particular seed $S$, translation gives a separate useful interface:
\[
 \card(t+S)=\card S,\qquad
 \sum_{x\in t+S}x=\card S\,t+\sum_{x\in S}x.
\]
When multiplication by $\card S$ is surjective, solve for $t$ to prescribe
the sum. Any translation-invariant property of $S$ is preserved.
\artifact{lean/SubsetTranslation.lean} proves this for arbitrary additive
commutative groups, then specializes to $\Zm T$ using coprimality.

This companion avoids the triangular formula and allows a noninterval
seed, but it does not remove coprimality for translates of that one seed.
The $T=4$, $S=\{0,1\}$ example explains exactly why. A physical-safety
property is not automatically translation-invariant; that premise must be
proved for the intended placement model. This companion and P1 target
overlapping interval machinery, so their savings should be counted once.

\section{Validation and reproducibility}
\label{sec:validation}

\subsection{What was actually checked}

The seven delivered proof experiments use Lean~4.32.0 and the repository-pinned mathlib
revision. They are ordinary kernel-checked Lean proofs. The final receipt
inventory distinguishes standalone generic theorems, a checked adapter
using freshly compiled repository prerequisites, and a focused production
module replacement.

\begin{longtable}{@{}p{0.39\textwidth}p{0.52\textwidth}@{}}
\toprule
Artifact & Checking scope \\
\midrule
\endhead
\artifact{RestrictedSubsetSums.lean} & Passed; explicit natural construction and all proper positive cyclic cardinalities, including noncoprime examples. \\
\artifact{AffineSubsetCoverage.lean} & Passed against four freshly compiled unchanged controller modules and the checked cyclic subset theorem. \\
\artifact{FixedPeriods.lean} & Passed; generic period transport and actual modular certificates, with an exact-order counterexample. \\
\artifact{ResidueTower.lean} & Passed; whole-expression transport and actual numeric specialization conditional on the three existing kernel-residue facts. \\
\artifact{PrimePowerAnnihilator.lean} & Passed; all-exponent divisibility, modular powers, and uniform exact order. \\
\artifact{LiftSeed.lean} & Passed against a freshly compiled baseline principal-unit module. \\
\artifact{SubsetTranslation.lean} & Passed; additive-group theorem, invariant-property variant, and cyclic specialization. \\
\bottomrule
\end{longtable}

The unchanged prerequisites \ident{ControllerIntervalSums},
\ident{ControllerOrbitTransport}, \ident{ControllerAuxiliaryChoice}, and
\ident{ControllerResidueLift} were compiled freshly for the affine adapter.
The uniform-order patch was checked separately under the real module name
\ident{DiophantineConsistency.PrimePowerOrder}; the unmodified
\ident{PrimePowerOrderChecks} module was then compiled against that fresh
replacement. The patch also passes \texttt{git apply --check} at the pinned
baseline. All six pre-existing public theorem signatures are preserved;
the only additional public theorem is \ident{pow_divisibility}.

The focused module changes from 80 to 79 physical lines, including the new
helper. This is the one measured source-size comparison in the deliverable.
Its principal benefit is a uniform theorem and fewer proof branches, not
a large line deletion. Compiler wall times are retained as receipts, but
the shared execution environment does not support a controlled timing
comparison.

\subsection{Trust and integration boundaries}

The 27 printed axiom reports from the delivered prototypes and the six
reports from the unchanged focused checks contain only the standard axioms
\ident{propext}, \ident{Classical.choice}, and \ident{Quot.sound}, or no
axioms. The package's validation script rejects failed compiler runs and
unexpected axiom dependencies. No supplied proof uses \ident{sorry},
\ident{admit}, a newly declared axiom, or \ident{native_decide}.

The explicit distinction between checking levels matters. A compiled
generic theorem with stated premises is not a rebuilt theorem about the
entire production kernel. The residue adapter accepts precisely the
already-proved kernel residue facts, while the affine adapter imports the
existing coordinate infrastructure and changes the endpoint cardinality
contract. The delivered patch has a narrower, genuinely replayed production
scope. No full-project build or complete consistency-equivalence proof is
claimed.

\subsection{Replay}

From the unpacked archive, use:
\begin{lstlisting}
python3 scripts/replay.py /path/to/ProveIt
pdflatex -interaction=nonstopmode -halt-on-error report.tex
pdflatex -interaction=nonstopmode -halt-on-error report.tex
\end{lstlisting}
The replay script verifies the repository pin and needed source hashes,
builds only the named prerequisites, compiles the proof experiments in an
isolated output directory, then checks the replacement module and unchanged
focused checks. It does not apply the patch or alter tracked source files.
The scripts and receipts separate reproducibility from any future decision
to adopt the changes.

The exclusion inventory can be regenerated without executing code from
the earlier archives:
\begin{lstlisting}
python3 scripts/build_exclusion_inventory.py \
  --repo /path/to/ProveIt --out regenerated-exclusion-audit
\end{lstlisting}


\section{Suggested implementation sequence}

\begin{enumerate}
\item Adopt the general cyclic subset-sum theorem and test its noncoprime
clients. Keep the old interval constructors temporarily if callers need
their explicit values.
\item Add the Boolean-free affine endpoint beside the old one. Make the
$a\le C_0$ capacity obligation visible at the exact physical composition
site. Only then retire the old arithmetic Boolean detour where appropriate.
\item Replace fixed-prime exact-order dependencies by period or projection
certificates. Split optional diagnostics into a separate module only if
reducing build dependencies is part of the goal.
\item Introduce the whole-coefficient residue theorem and instantiate it
with the three existing literal-kernel residue facts. Preserve $r_w$ and
$n$ as independent parameters and retain their zero cases.
\item Apply the small uniform-annihilator patch and its focused checks.
Choose the first-lift seed adapter when it simplifies the current-generator
interface enough to justify the helper's maintenance cost.
\end{enumerate}

For each adoption, measure source size after including helpers and retained
compatibility declarations. Separately record build time, memory, and axiom
dependencies. A smaller source file is not evidence of faster elaboration;
an improved theorem interface can be worthwhile even when its first
implementation has approximately the same line count.

\section{Source and artifact guide}

The claims about current code are grounded in the pinned source links
throughout the report. The exclusion locations are
\href{\repobase/Computability/DiophantineConsistency/research/lean-proof-simplification/README.md}{the placed-packet index},
\href{\repobase/Computability/DiophantineConsistency/research/lean-proof-simplification/ADOPTION.md}{the adoption record}, and
\href{\repobase/docs/incoming/README.md}{the incoming index}.
Standard library support is from the pinned
\href{https://github.com/leanprover-community/mathlib4/blob/81a5d257c8e410db227a6665ed08f64fea08e997/Mathlib/GroupTheory/OrderOfElement.lean}{multiplicative-order development},
\href{https://github.com/leanprover-community/mathlib4/blob/81a5d257c8e410db227a6665ed08f64fea08e997/Mathlib/NumberTheory/Multiplicity.lean}{LTE and multiplicity development}, and
\href{https://github.com/leanprover-community/mathlib4/blob/81a5d257c8e410db227a6665ed08f64fea08e997/Mathlib/Data/Nat/PrimeFin.lean}{finite prime-factor interface}.

The archive contains the complete LaTeX source and PDF, Lean experiments,
the focused patch, proof-check receipts, a replay script, source and exclusion
inventories, and independent mathematical-review notes. It is a reviewable
research deliverable; none of these proposals was pushed to the repository.

\end{document}
